\section*{अध्याय \ref{ch:b3:bioinorganic} --- जैवअकार्बनिक रसायन}

\begin{solution}{exo:b3:bioinorganic:1}
$\log(\theta/(1 - \theta))$, $\log 0.25 = -0.602$ से $\log 4 = +0.602$ तक जाता है, जबकि $\log p$
$\log(5.9/2.1) = 0.449$ बढ़ता है: $n = 1.204/0.449 = 2.7$।
\end{solution}

\begin{solution}{exo:b3:bioinorganic:2}
$\ce{Mn^2+} < \ce{Fe^2+} < \ce{Co^2+} < \ce{Ni^2+} < \ce{Cu^2+} > \ce{Zn^2+}$। ज़िंक ($d^{10}$) में न लिगैंड-क्षेत्र स्थायीकरण है और न
यान--टेलर विकृति: यह कॉपर से नीचे लौट आता है।
\end{solution}

\begin{solution}{exo:b3:bioinorganic:3}
$\ce{Ca^2+}$ और $\ce{Fe^3+}$, कठोर: कार्बोक्सिलेट (और आयरन के लिए फ़ीनॉलेट)। $\ce{Cu+}$ और $\ce{Hg^2+}$, मृदु:
सिस्टीन थायोलेट। $\ce{Zn^2+}$, सीमांत: हिस्टिडीन और सिस्टीन।
\end{solution}

\begin{solution}{exo:b3:bioinorganic:4}
$\ce{2H2O -> O2 + 4H+ + 4e-}$: प्रति $\ce{O2}$ चार इलेक्ट्रॉन, चार फ़ोटॉन; प्रति मोल $\ce{O2}$
फ़ोटॉनों के चार मोल।
\end{solution}

\begin{solution}{exo:b3:bioinorganic:5}
मायोग्लोबिन: $\theta = 13/13.37 = 0.972$ और $5/5.37 = 0.931$: यह अपनी ऑक्सीजन का $(0.972 - 0.931)/0.972 = 4$\,\%
छोड़ता है। हीमोग्लोबिन: $\theta = 0.972$ और $0.724$: यह 26\,\% छोड़ता है।
\end{solution}

\begin{solution}{exo:b3:bioinorganic:6}
$[\mathrm{HbCO}]/[\mathrm{HbO_2}] = 230 \times \num{50e-6}/0.2095 = 0.055$; अंश $0.055/1.055 = 5.2$\,\%।
\end{solution}

\begin{solution}{exo:b3:bioinorganic:7}
सल्फ़ाइड प्रत्येक $-2$, सिस्टीनेट प्रत्येक $-1$। $\ce{[Fe2S2(SR)4]^2-}$: आयरनों का योग $-2 + 4 + 4 = +6$: दो
Fe(III)। $\ce{[Fe2S2(SR)4]^3-}$: $+5$, एक Fe(III) और एक Fe(II) (मिश्रित संयोजकता)। $\ce{[Fe4S4(SR)4]^2-}$:
$-2 + 8 + 4 = +10$, दो Fe(III) और दो Fe(II), प्रत्येक आयरन औसतन $+2.5$।
\end{solution}

\begin{solution}{exo:b3:bioinorganic:8}
$\ce{[PtCl2(NH3)2] + H2O -> [PtCl(H2O)(NH3)2]+ + Cl-}$। रक्त की उच्च क्लोराइड सांद्रता
साम्य को वापस धकेलती है; कोशिकाओं के भीतर क्लोराइड कहीं कम है, और ऐक्वा
संकुल बनता है। प्लैटिनम ग्वानीन के N7 से बँधता है। \textit{trans} समावयवी में दोनों
\omterm{def:b3:complex-mechanisms:labile}{विनिमयशील} स्थल एक-दूसरे के सम्मुख हैं और एक रज्जुक के दो संलग्न क्षारकों तक
नहीं पहुँच सकते।
\end{solution}

\begin{solution}{exo:b3:bioinorganic:9}
$1/T_1 = 1/\qty{1.2}{s} + \qty{4.0}{L.mmol^{-1}.s^{-1}} \times \qty{0.10}{mmol/L} = 0.833 + 0.400 = \qty{1.23}{s^{-1}}$: $T_1 = \qty{0.81}{s}$।
\end{solution}

\begin{solution}{exo:b3:bioinorganic:10}
$K = K_{\mathrm{PbL}}/K_{\mathrm{CaL}} = 10^{18.0 - 10.7} = 10^{7.3} = \num{2e7}$: सीसा कैल्सियम को पूरी तरह विस्थापित कर देता है। मुक्त
लिगैंड रक्त के कैल्सियम को भी बाँध लेता और उसकी सांद्रता को ख़तरनाक रूप से
घटा देता; कैल्सियम संकुल सीसे से विनिमय करता है और कैल्सियम
स्तर अपरिवर्तित छोड़ता है।
\end{solution}

\begin{solution}{exo:b3:bioinorganic:11}
$K_M$ से बहुत नीचे अभिक्रिया प्रथम कोटि की है: $k' = (k_{\mathrm{cat}}/K_M)[\mathrm E]$, $t_{1/2} = \ln 2/k'$। कार्बोनिक
ऐनहाइड्रेज़: $k' = \qty{100}{s^{-1}}$, $t_{1/2} = \qty{6.9}{ms}$। विशिष्ट \omterm{def:b3:complex-kinetics:enzyme}{एंज़ाइम}: $k' = \qty{0.10}{s^{-1}}$, $t_{1/2} = \qty{6.9}{s}$, एक हज़ार
गुना अधिक, उस एक सेकंड के लिए बहुत धीमा जो एक लाल कोशिका फेफड़े की केशिका में बिताती है।
\end{solution}

\begin{solution}{exo:b3:bioinorganic:12}
आधे स्थलों पर CO का अर्थ है $[\mathrm{HbCO}] = [\mathrm{HbO_2}]$: $p(\ce{CO}) = p(\ce{O2})/M = 0.2095/230 = \qty{9.1e-4}{bar}$, लगभग
\qty{910}{ppm}।
\end{solution}

\begin{solution}{pb:b3:bioinorganic:1}
\textbf{1.} $\theta = 13^{2.7}/(3.5^{2.7} + 13^{2.7}) = 0.972$।

\textbf{2.} $\theta(\qty{5.0}{kPa}) = 0.724$।

\textbf{3.} $5.0/(0.37 + 5.0) = 0.931$।

\textbf{4.} इसका $p_{50}$ दस गुना कम है: कार्यरत पेशी के दाबों पर यह उसी दाब पर हीमोग्लोबिन से कहीं अधिक
संतृप्त होता है, अतः ऑक्सीजन
हीमोग्लोबिन से मायोग्लोबिन की ओर जाती है।

\textbf{5.} $n = 1$।

\textbf{6.} एक \omterm{def:b3:bioinorganic:haem}{हीम} पर बंधन उसके आयरन को तल में ले आता है और हिस्टिडीन तथा
उसकी कुंडलिनी को खींचता है; उप-इकाइयाँ मिलकर एक ऐसे रूप में खिसकती हैं जिसकी
अन्य स्थलों के लिए बंधुता अधिक है।

\textbf{7.} $\qty{150}{g/L}/\qty{64500}{g/mol} = \qty{2.33e-3}{mol/L}$।

\textbf{8.} $4 \times \num{2.33e-3} = \qty{9.30e-3}{mol/L}$।

\textbf{9.} $\qty{9.30}{mmol/L} \times 0.972 = \qty{9.04}{mmol}$।

\textbf{10.} $\qty{9.30}{mmol/L} \times (0.972 - 0.724) = \qty{2.31}{mmol}$।

\textbf{11.} $\qty{2.31e-3}{mol} \times \qty{22.41}{L/mol} = \qty{52}{mL}$।

\textbf{12.} $0.248/0.972 = 26$\,\%।

\textbf{13.} $5.0^{2.7}/(4.0^{2.7} + 5.0^{2.7}) = 0.646$।

\textbf{14.} $\qty{9.30}{mmol/L} \times (0.972 - 0.646) = \qty{3.03}{mmol}$।

\textbf{15.} $3.03/2.31 = 1.31$: तीस प्रतिशत अधिक।

\textbf{16.} श्वसन से बनी कार्बन डाइऑक्साइड, जिसे कार्बोनिक ऐनहाइड्रेज़
कार्बोनिक अम्ल में जलयोजित करता है, और कठोर श्रम में लैक्टिक अम्ल।

\textbf{17.} $\qty{3.03}{mmol/L} \times \qty{5.0}{L/min} = \qty{15.1}{mmol/min}$।

\textbf{18.} $\qty{15.1e-3}{mol/min} \times \qty{22.41}{L/mol} = \qty{0.34}{L/min}$।

\textbf{19.} $230 \times \num{100e-6}/0.2095 = 0.110$।

\textbf{20.} $0.110/1.110 = 0.099$: स्थलों का लगभग 10\,\%।

\textbf{21.} जो स्थल अब भी ऑक्सीजन वहन करते हैं, CO भी वहन करने वाले अणु में, वे
उच्च-बंधुता रूप में होते हैं: वक्र बाईं ओर खिसकता है और वे ऊतकों में अपनी
ऑक्सीजन कम छोड़ते हैं।

\textbf{22.} स्पर्धा संबंध से अनुपात $[\mathrm{HbCO}]/[\mathrm{HbO_2}]$, $p(\ce{O2})$ के अनुपात में गिरता है:
शुद्ध ऑक्सीजन में साँस लेना, वायु के दाब का लगभग पाँच गुना, कार्बन
मोनोऑक्साइड को तेज़ी से विस्थापित करता है।

\textbf{23.} CO के रैखिक बंधन में बाधा डालकर और $\ce{O2}$ से हाइड्रोजन आबंध बनाकर, यह $M$ को
मुक्त \omterm{def:b3:bioinorganic:haem}{हीम} के कहीं बड़े मान के बजाय सैकड़ों में रखता है; अन्यथा शरीर में बनी
कार्बन मोनोऑक्साइड भी हीमोग्लोबिन के बड़े भाग को अवरुद्ध कर देती।

\textbf{24.} विश्राम पर रक्त प्रति मिनट लगभग \qty{0.34}{L} ऑक्सीजन पहुँचाता है।
\end{solution}
